% 参考资料章节

\clearpage
\phantomsection
\label{sec:refs}
\renewcommand{\refname}{参考资料}
\addcontentsline{toc}{section}{参考资料}

\begin{thebibliography}{99}

\bibitem{ref1}
Liu, C. L., Layland, J. W. Scheduling algorithms for multiprogramming in a hard-real-time environment.
\emph{Journal of the ACM}, 20(1): 46--61, 1973.

\bibitem{ref2}
Davis, R. I., Burns, A. A survey of hard real-time scheduling for multiprocessor systems.
\emph{ACM Computing Surveys}, 43(4): 1--44, 2011.

\bibitem{ref3}
Bastoni, A., Brandenburg, B. B., Anderson, J. H. Cache-related preemption and migration delays:
empirical approximation and impact on schedulability.
In \emph{Proc. OSPERT}, 2010.

\bibitem{ref4}
Lunniss, W., Altmeyer, S., Davis, R. I. Optimising task assignment in real-time multiprocessor systems.
\emph{Real-Time Systems}, 49(4): 465--496, 2013.

\bibitem{ref5}
Hennessy, J. L., Patterson, D. A. \emph{Computer Architecture: A Quantitative Approach}, 6th ed.
Morgan Kaufmann, 2019.

\bibitem{ref6}
Ding, C., Zhong, Y. Predicting whole-program locality through reuse distance analysis.
In \emph{Proc. PLDI}, 2003.

\bibitem{ref7}
Zhao, S., Dai, X., Lesage, B., Bate, I. Cache-aware allocation of parallel jobs on multi-cores
based on learned recency.
In \emph{Proc. 31st International Conference on Real-Time Networks and Systems (RTNS)}, 2023.

\bibitem{ref8}
Buttazzo, G. C. \emph{Hard Real-Time Computing Systems: Predictable Scheduling Algorithms and
Applications}, 3rd ed. Springer, 2011.

\bibitem{ref9}
Linux kernel source tree: \texttt{kernel/sched/core.c}, \texttt{include/linux/sched.h},
\texttt{Documentation/scheduler/}. Version 6.8.

\bibitem{ref10}
Intel Corporation. \emph{Intel 64 and IA-32 Architectures Optimization Reference Manual}, 2021.

\end{thebibliography}